\appendix

\chapter{Függelékek}
\label{chap:fuggelek}

A függelék az „Automatikus hangszerfelismerés zenefelvételek alapján” című
államvizsga-szakdolgozathoz tartozó digitális projektanyag szerkezetét foglalja
össze. A cél nem a teljes forráskód bemásolása, hanem annak áttekinthető
megmutatása, hogy a nyers hangminták, a kísérleti adathalmazok, a kinyert
jellemzők, a betanított modellek és a futtató szkriptek hol találhatók.

%% ═══════════════════════════════════════════════════════════════════════════
\section{A projekt könyvtárszerkezete}
\label{sec:konyvtarszerkezet}

A benyújtott digitális függelék az alábbi főbb elemeket tartalmazza. A
forráskód és a dokumentáció a GitHubon elérhető nyilvános Git
adattárban\footnote{\url{https://github.com/nCsab/Instrument-Recognition}} is
megtalálható. A nagy méretű, helyben gyűjtött vagy generált adathalmazok nem
kerülnek nyilvános terjesztésre; a leadott forráscsomag a futtatáshoz szükséges
szkripteket, a dolgozat forrását és a kiválasztott valós idejű checkpointot
tartalmazza. A számozott szkriptek a feldolgozási csővezeték egymásra épülő
lépéseit jelzik.

\begin{table}[ht!]
\centering
\small
\caption{A projekt legfelső szintű könyvtárszerkezete}
\label{tab:konyvtarszerkezet}
\begin{tabular}{p{4.5cm}p{10cm}}
\toprule
\textbf{Könyvtár / Fájl} & \textbf{Leírás} \\
\midrule
\texttt{thesis/} & A \LaTeX\ formátumú szakdolgozat forrásfájljai, stílusfájl, irodalomjegyzék, ábrák és a generált \texttt{main.pdf}. \\
\texttt{scripts/} & A teljes feldolgozási, jellemzőkinyerési, tanítási és valós idejű felismerési folyamat Python szkriptjei. \\
\texttt{owndataset/dataset.py} & Manuális segédszkript hosszabb hangfájlokból 5~másodperces blokkok kivágásához. A nyers saját hangfájlok nem részei a nyilvános csomagnak. \\
\texttt{models/exp\_final/} & A valós idejű demonstrációhoz használt kiválasztott Log-Mel + 2D CNN checkpoint. A teljes kísérleti modell- és riportmappa helyben maradt. \\
\texttt{README.md} & Rövid futtatási leírás, pipeline-összefoglaló és adatkezelési megjegyzések. \\
\texttt{requirements.txt} & A lokális futtatáshoz rögzített Python-csomagverziók. \\
\bottomrule
\end{tabular}
\end{table}
\clearpage
%% ═══════════════════════════════════════════════════════════════════════════
\section{A feldolgozási csővezeték szkriptjei}
\label{sec:pipeline_scripts}

A rendszer teljes feldolgozási lánca indexelt Python szkriptekből áll, amelyek
a tervezés sorrendjében futtathatók. Az alábbi táblázat az aktuális, végleges
pipeline főbb elemeit foglalja össze.

\begin{table}[ht!]
\centering
\small
\caption{A feldolgozási csővezeték szkriptjei és feladataik}
\label{tab:pipeline_scripts}
\begin{tabular}{p{5.8cm}p{8.7cm}}
\toprule
\textbf{Szkript} & \textbf{Feladat} \\
\midrule
\texttt{01\_prepare\_\allowbreak clean\_\allowbreak dataset.py} & A saját clean adathalmaz elkészítése: forrásblokkok szeparált felosztása, 1~másodperces szeletelés és zaj/csend minták kezelése. \\
\texttt{02\_acquire\_\allowbreak mic\_\allowbreak data.py} & A clean minták visszajátszásának, mikrofonos rögzítésének és split/osztály szerinti rendezésének támogatása. \\
\texttt{03a\_merge\_\allowbreak mic\_\allowbreak to\_\allowbreak train.py} & A négy kísérleti adathalmaz összeállítása a clean, augmentált és mikrofonos mintákból. \\
\texttt{03b\_check\_\allowbreak dataset\_\allowbreak counts.py} & A kísérleti adathalmazok osztályonkénti és split szerinti ellenőrzése. \\
\texttt{04a\_extract\_\allowbreak features.py} & A jellemzők kinyerése és normalizálása. A clean baseline esetén MFCC, STFT és Log-Mel, a további kísérleteknél Log-Mel, az \texttt{exp\_final} esetén pedig YAMNethez nyers audio is készül. \\
\texttt{04b\_check\_\allowbreak extracted\_\allowbreak counts.py} & A kinyert NumPy tömbök és címkék méretének ellenőrzése. \\
\texttt{05a\_realtime\_\allowbreak recognition.py} & A végső valós idejű Log-Mel + 2D CNN felismerő alkalmazás. \\
\texttt{06\_train\_\allowbreak model\_\allowbreak colab/06\_train\_\allowbreak custom\_\allowbreak cnn.py} & A saját 2D CNN modell Google Colab környezetben futtatott tanítása, checkpointolása és riportgenerálása. \\
\texttt{06\_train\_\allowbreak model\_\allowbreak colab/06\_train\_\allowbreak yamnet.py} & A YAMNet embeddingekre épülő transfer learning referencia tanítása és kiértékelése. \\
\texttt{scripts/utils/} & Segédfüggvények az augmentációhoz és a jellemzőkinyeréshez. \\
\bottomrule
\end{tabular}
\end{table}

\clearpage
%% ═══════════════════════════════════════════════════════════════════════════
\section{Fejlesztői és futtatási környezet}
\label{sec:fejlesztoi_kornyezet}

Az alábbi táblázat tartalmazza a projekt fejlesztéséhez és futtatásához szükséges szoftverfüggőségeket.

\begin{table}[H]
\centering
\caption{A fejlesztői környezet főbb összetevői}
\label{tab:fejlesztoi_kornyezet}
\begin{tabular}{p{4.2cm}p{2.5cm}p{7.5cm}}
\toprule
\textbf{Szoftver / Könyvtár} & \textbf{Verzió} & \textbf{Felhasználás} \\
\midrule
Python & 3.10+ & A teljes feldolgozási csővezeték nyelve \\
\addlinespace
TensorFlow & 2.19.1 & A 2D CNN modell futtatása és mentése \\
\addlinespace
Keras & 3.13.2 & Neurális hálózati rétegek és modellkezelés \\
\addlinespace
Librosa & 0.11.0 & Hangfájlok betöltése, STFT, mel-spektrogram és MFCC kinyerés \\
\addlinespace
NumPy & 2.1.3 & Jellemzőmátrixok tárolása és manipulálása \\
\addlinespace
SoundFile & 0.13.1 & Hangfájlok olvasása és írása \\
\addlinespace
SoundDevice & 0.5.2 & Valós idejű mikrofon-hozzáférés \\
\addlinespace
scikit-learn & 1.6.1 & Classification report és konfúziós mátrix előállítása \\
\addlinespace
Matplotlib & 3.10.3 & Diagramok, spektrogramok és kiértékelési ábrák készítése \\
\addlinespace
Seaborn & 0.13.2 & Konfúziós mátrixok és F1-score ábrák formázása \\
\addlinespace
Google Colab & -- & GPU-gyorsított modelltanítás felhőkörnyezetben \\
\addlinespace
Git & 2.x & Verziókezelés és a projekt előrehaladásának dokumentálása \\
\addlinespace
\LaTeX\ (TeX Live 2026) & -- & A szakdolgozat szedése a Sapientia sablon alapján \\
\bottomrule
\end{tabular}
\end{table}
